\phantomsection
\addcontentsline{toc}{chapter}{ABSTRACT}
\begin{center}{\bfseries\fontsize{16}{20}\selectfont ABSTRACT}\end{center}
\vspace{6pt}
This project designs, implements and evaluates machine-learning pipelines across all five units of the course, turning
raw data into reliable and interpretable decisions. Its objectives are to understand data through exploratory analysis,
to build and compare supervised models, to discover structure without labels, to model sequential behaviour, and to
investigate the effectiveness of ensemble classifiers. In PBL 1, the Google Play Store and Video Game Sales datasets
were cleaned and explored, and the Combined Cycle Power Plant and Bank Marketing datasets were used for regression and
classification, where tuned XGBoost reached an RMSE of 3.17 MW (R² 0.966) and an imbalance-aware XGBoost classifier
reached an F1-score of 0.531 with a 3.3-fold lift. In PBL 2, the UCI Human Activity Recognition dataset was used:
K-Means and Gaussian Mixture clustering recovered activity structure with up to 61.2 percent accuracy without labels,
and Hidden Markov Models with hand-implemented forward–backward and Viterbi algorithms decoded the activity sequences
of unseen people with 98.7 percent accuracy. In PBL 3, ten ensemble classifiers, spanning bagging, boosting, voting and
stacking, were compared with Logistic Regression and Decision Tree baselines on \dbRowsClean{} patient records for
diabetes prediction. The best model, \dbBestName{} with an optimised decision threshold, achieved an F1-score of
\dbBestFone{} and a ROC-AUC of \dbBestAuc{} on unseen patients, significantly better than Logistic Regression (F1
\dbLrFone). Across the project, careful leakage control, metrics suited to the class balance and the combination of
multiple learners consistently produced the most accurate and trustworthy models.
